\subsection{A uniform category of coefficient groups}
Theorem~\ref{thm:uniform-coefficient-category} replaces the
family-dependent coefficient choice by a single category,
retaining the ordinary ideal-generated projective problem.
We use the standard covariant module convention for a
preadditive category; representables, finite projectors and
extension along an additive functor have their usual meanings
\cite{Mitchell}. We use the covariant convention and representable
projectivity described in \cite[\S1, Proposition~3, pp.~1--4]{Murfet};
extension and its adjunctions are given in
\cite[\S4, Theorem~18 and Proposition~19, pp.~23--25]{Murfet}.
The coefficient category and its ideal below specify the form
needed for inverse-semiring amplitudes.

\begin{proof}[Proof of Theorem~\ref{thm:uniform-coefficient-category}]
Admissibility and $Q$-linearity give, in the required order,
\[
 q_{zy}q_{yx}
 =\phi_x\bigl(\phi_y(l_z)l_y\bigr)
 \leq\phi_x(l_z)=q_{zx}.
\]
Hence $q_{xx}^2\leq q_{xx}$ and $1\vee q_{xx}$ is an upper
multiplicative idempotent. Accounting for a diagonal identity
term also gives
\[
 d_{zy}d_{yx}\leq d_{zx},\qquad
 d_{zy}q_{yx}\leq q_{zx},\qquad
 q_{zy}d_{yx}\leq q_{zx}.
\]
For instance the extra term in the first factor occurs only
when $z=y$, and is then $q_{yx}=q_{zx}$.
These inequalities define composition in $\mathcal C_L$.
The coefficient identities, associativity and naturality give
its unital category axioms. Every $\mathcal N_L$-arrow has a
representative at its $q$-support, so the two mixed inequalities
prove both ideal closures. The diagonal quotient follows.
The set of objects is a subset of $L\times Q^L$, hence is small.

Since $d_{yx}l_x\leq l_y$, every coefficient defines an action
arrow, giving the unital functor $\pi$. Its image on
$\mathcal N_L$ lies in $\mathcal J_L$ by admissibility.
Write $K=\ker\pi$. Then
\begin{equation}\label{eq:uniform-annihilator}
 K\mathcal N_L=0.
\end{equation}
Indeed, if $a\in K(y,z)$, the directed-colimit zero criterion
for $\mathcal A_L(y,z)$ supplies $r\geq d_{zy}$ with
$rl_y\leq l_z$ at which $a$ becomes zero. For an
$\mathcal N_L(x,y)$-arrow represented at $q_{yx}$,
\[
 rq_{yx}=\phi_x(rl_y)\leq\phi_x(l_z)=q_{zx}.
\]
Compute the product by first transporting $a$ to that zero
and then transporting the product to $d_{zx}$. This proves
\eqref{eq:uniform-annihilator}. No reverse annihilator,
injectivity of transitions, or vanishing of $F(0)$ is used.

For finite lists of objects, let $e,f$ be idempotent matrices
with entries in $\mathcal N_L$. Matrix calculus uses the common
opposite convention for covariant representables on both sides.
All mixed action corners
lift exactly to coefficient corners. To see surjectivity,
represent an intervening action arrow $b_{jk}$ at $r$ with
$rl_k\leq l_j$. The two outside entries have representatives
at $q_{ij}$ and $q_{km}$, and
\[
 rq_{km}=\phi_m(rl_k)\leq q_{jm},\qquad
 q_{ij}rq_{km}\leq q_{ij}q_{jm}\leq q_{im}.
\]
Thus each term of $\pi(e)b\pi(f)$ has a representative in
$F(q_{im})$, hence in $\mathcal N_L$. Choose a coefficient
matrix lift $c$ and replace it by $ecf$. If a corner element
$z\in e\operatorname{Mat}(\mathcal C_L)f$ maps to zero,
then $z=zf=0$ by \eqref{eq:uniform-annihilator}.
Consequently
\[
 e\operatorname{Mat}(\mathcal C_L)f
 \xrightarrow{\sim}
 \pi(e)\operatorname{Mat}(\mathcal A_L)\pi(f)
\]
is an additive isomorphism compatible with composition.
Object repetitions in these finite lists are allowed.
A matrix entry between two copies of the same categorical
object belongs to $F(1\vee q_{xx})$, and its ideal entry is
still represented at $q_{xx}$; the calculations remain valid.

A finite projective $P$ over a preadditive category satisfies
$\mathcal NP=P$ exactly when it has an idempotent presentation
with entries in $\mathcal N$. For a split presentation
$P\to V\to P$ by finite representables, the section sends
$\mathcal NP$ into $\mathcal NV$, proving necessity.
Conversely an $\mathcal N$-entry projector $e$ satisfies
$eV=e^2V\subseteq e(\mathcal NV)=\mathcal N(eV)$.
Left Kan extension sends the representable at $x$ to the
representable at $l_x$, and therefore sends finite projectors
to their images under $\pi$. The exact corner calculation
proves full faithfulness on the stated subcategory.

For essential surjectivity, present a finite projective
$\mathcal J_L$-generated amplitude by an idempotent $E$ with
entries in $\mathcal J_L$. In each source column $j$, choose
admissible functionals representing its finitely many entries
and replace them by their finite pointwise join $\phi_j$.
Every entry then transports to $F(\phi_j(l_i))$.
Choose the objects $(l_j,\phi_j)$ of $\mathcal C_L$.
The matrix $E$ lifts to a matrix $a$ with entries in
$\mathcal N_L$, even when some of these objects coincide.
Its defect belongs to $K$, so
\[
 (a^2-a)a=0,\qquad a^3=a^2,\qquad (a^2)^2=a^2.
\]
Thus $a^2$ is an exact coefficient projector lifting $E$.
This proves essential surjectivity and the equivalence.
Extension to the quotient category is $P\mapsto
P/\mathcal N_LP$, proving the asserted kernel description.

Finally, $u\in\operatorname{Aut}_Q(L)$ sends an admissible
pair to $(u(l),\phi u^{-1})$, since
\[
 \phi(u^{-1}(k))u(l)=u\bigl(\phi(u^{-1}(k))l\bigr)\leq k.
\]
All $q$-values are unchanged, so this is a relabeling of the
same coefficient groups and products. It commutes with $\pi$
and the action-category relabeling. The canonical projectivity
criterion now gives the fixed-support classification and the
automorphism-orbit statement.
\end{proof}

The coefficient category may still be infinite. In the
singleton-support arbitrary-ideal example of
Proposition~\ref{prop:canonical-arbitrary-ideal}, it is the
single ring $R$ and its ideal is the arbitrary ordinary ideal
$I$. Thus the theorem removes the action-arrow colimits and
family-dependent coefficient choices, while retaining the
ordinary ideal-generated projective problem. It makes no
finite-object assertion and does not contradict
Example~\ref{ex:no-fixed-finite-objects}.
